In summary, we demonstrated how to work with topological representations
of various \textbf{PPAD} problem, and how under certain conditions
one can parsimoniously reduce them to \textit{PureCircuit}.
First we showed how \textit{Kleene logic} can be embedded
into \textbf{PPAD} using as a primary example the \textit{StrongSperner}
problem and how certain problems
retain their \textbf{PPAD}-completeness. Morevover, we demonstrated that under hazard-free conditions
and discreteness, how
one can parsimoniously reduce \textit{HF-DiscreteStrongSperner} to
\textit{PureCricuit}.  
We were also able to construct a parsimonious reduction
from the \textit{EndOfLine} problem to  \textit{StrongSperner} instrances
with discrete and antipodal properties. 
Since we know that the hazard free version of \textit{AntipodalStrongSperner} 
can be parsimoniously reduced to the \textit{PureCircuit}, we argue
that if one can construct a hazard-free version of the circuit,
we can achieve the full parsimonious reduction chain between \textit{EndOfLine}.


Extending this idea, we showed how \textit{SourceOrExcess}$(2,1)$ can also
be encoded topologically and demonstrated how 
the dimensionality of the space we encode it can affect the counting hardness
of problems. We argue that this has some substantial consequences as first
we found another \textbf{PPAD}-complete problem with no
clear combintorial interpretation. Moreover, this
lead to interesting questions as one can ask how does the degree of a graph
affect the smallest possible space we can encode it in.

Moreover we investigate the different counting classes that
can be represented in \textbf{PPAD} by expanding on Hollender et al.'s work \cite{hollender_ComplexityMultisourceVariants_2018}.
Specifically we expanded on the combinatorial argument their combinatorial to
argued about the infinite hiararchy of \textbf{PPAD} problems and their minimum number of solutions.
In addtion, we argued the importance of the \textit{Unbalanced} problem as an
important stepping stone to a total \textit{PureCircuit} Cook-Levin theorem,
due to its overall relation with all the other unbalanced digraph problems.


Lastly, we showcased our work on our \textit{PureCircuit} visualisation software.
We demonstrated its core functionalities such as \textit{PureCircuit} editor
tool with assignment verification.  Moreover we implemented several
metaheuristic algorithms for assignment finding assignments in large graphs, such
as \textit{Hill Climbing} and \textit{Genetic Algorithms}. Most importantly,
we also implemented a specialised backtracking algorithm whick takes advantage
of the structure of the graph. Despite its high worst-case running time,
we manage to use several greedy methods to cut down breadth and with of the tree
making it ideal for verification of small circuits.






